\section*{अध्याय \ref{ch:b3:atmospheric-dynamics} --- वायुमंडलीय गतिकी}

\begin{solution}{exo:b3:atmospheric-dynamics:1}
(क)~$\zeta = 2\omega = 2\times2\pi/86\,400 = \qty{1.45e-4}{s^{-1}}$।
(ख)~$f = \num{7.29e-5}$ और \qty{1.26e-4}{s^{-1}}: \omterm{def:b3:atmospheric-dynamics:eddy}{भँवर} की \omterm{def:b3:atmospheric-dynamics:vorticity-kinds}{ग्रहीय भ्रमिलता} लगभग
\ang{30} N की दोगुनी है और
\ang{60} N की उससे कुछ अधिक।
\end{solution}

\begin{solution}{exo:b3:atmospheric-dynamics:2}
$\zeta = f_{30} - f_{40} = \num{7.29e-5} - \num{9.37e-5} =
\qty{-2.08e-5}{s^{-1}}$: प्रतिचक्रवातीय।
\end{solution}

\begin{solution}{exo:b3:atmospheric-dynamics:3}
$k = 2\pi/(5\times10^6) = \qty{1.257e-6}{m^{-1}}$; $c = -\beta/k^2 =
-\num{1.62e-11}/\num{1.58e-12} = \qty{-10.3}{m/s}$: पश्चिम की ओर।
\end{solution}

\begin{solution}{exo:b3:atmospheric-dynamics:4}
$\beta = 2\Omega\cos\ang{60}/a = \qty{1.14e-11}{m^{-1}.s^{-1}}$;
$\lambda_s = 2\pi\sqrt{10/\num{1.14e-11}} = \qty{5870}{km}$; वृत्त,
\qty{20015}{km} में $n_s = 3.4$ तरंगें समाती हैं।
\end{solution}

\begin{solution}{exo:b3:atmospheric-dynamics:5}
$q = f/h$ संरक्षित रहता है: परास के ऊपर $\zeta = f(8/10) - f =
\qty{-2.06e-5}{s^{-1}}$, प्रतिचक्रवातीय; बेसिन के ऊपर $f(12/10) - f =
\qty{+2.06e-5}{s^{-1}}$, चक्रवातीय। जहाँ स्तंभ खिंचते हैं, परास की पीठ-पीछे की ओर
और बेसिन के ऊपर, वहाँ चक्रवातीय भ्रमिलता उत्पन्न होती है:
वहाँ अवदाब बनने की प्रवृत्ति होती है।
\end{solution}

\begin{solution}{exo:b3:atmospheric-dynamics:6}
$k = 5/(a\cos\ang{45}) = \qty{1.11e-6}{m^{-1}}$, $\beta/k^2 =
\qty{13.1}{m/s}$: $c = \qty{1.9}{m/s}$, $c_g = \qty{28.1}{m/s}$। ऊर्जा
\num{2.1}~दिनों में \qty{5000}{km} तक फैलती है, प्रतिरूप \num{31}~दिनों में।
\end{solution}

\begin{solution}{exo:b3:atmospheric-dynamics:7}
(क)~$k^2 = \beta/U - l^2 = \num{1.079e-12} - \num{2.74e-13} =
\qty{8.05e-13}{m^{-2}}$, $k = \qty{8.97e-7}{m^{-1}}$: $\lambda =
\qty{7000}{km}$ और $n = 4.0$, $l = 0$ की तुलना में कम तरंगें। (ख)~$l^2 =
\num{1.097e-12} > \beta/U = \num{1.079e-12}$: $k^2$ ऋणात्मक होगा।
\end{solution}

\begin{solution}{exo:b3:atmospheric-dynamics:8}
$L_d = 0.012\times10^4/\num{1.031e-4} = \qty{1164}{km}$; $\sigma =
0.31\times\num{1.031e-4}\times\num{2e-3}/0.012 = \qty{5.3e-6}{s^{-1}}$, एक
$e$-गुना समय \num{2.2}~दिन; तरंगदैर्घ्य $2\pi\times1164/1.61 =
\qty{4540}{km}$।
\end{solution}

\begin{solution}{exo:b3:atmospheric-dynamics:9}
$[v^*T^*] = (VT_0/2)\cos\delta$: \qty{25}{K.m/s}, \qty{12.5}{K.m/s} और 0;
ऊष्मा फ़्लक्स \num{25100}, \num{12600} और \qty{0}{W/m^2}। जब तापमान
पवन से एक चौथाई तरंगदैर्घ्य कला से बाहर होता है, तो उतनी ही गर्म हवा दक्षिण
को जाती है जितनी उत्तर को: कोई परिवहन नहीं।
\end{solution}

\begin{solution}{exo:b3:atmospheric-dynamics:10}
$U$ के सापेक्ष रैखिकीकरण, जहाँ $\zeta' = \nabla^2\psi'$ और $v' =
\partial_x\psi'$: $(\partial_t + U\partial_x)\nabla^2\psi' + \beta\partial_x
\psi' = 0$। $\psi' \propto \mathrm e^{\mathrm i(kx + ly - \omega t)}$ के
साथ, $(\omega - Uk)K^2 = -\beta k$, अतः $\omega = Uk - \beta k/K^2$। वायु के
सापेक्ष कला-वेग प्रत्येक तरंग के लिए $c - U = -\beta/K^2 < 0$ है।
\end{solution}

\begin{solution}{exo:b3:atmospheric-dynamics:11}
$\varphi_H = [5\times9.81\times12\,000\times0.25/(3\times(\num{7.292e-5})^2
\times(\num{6.371e6})^2)]^{1/2} = \qty{0.477}{rad}$, \ang{27.3}।
$u_M = \Omega a\varphi_H^2 = \qty{106}{m/s}$; सटीक रूप में, $\Omega a\sin^2
\varphi_H/\cos\varphi_H = \qty{110}{m/s}$: प्रेक्षित जेट से
ढाई गुना।
\end{solution}

\begin{solution}{exo:b3:atmospheric-dynamics:12}
(क)~पट्टी के आर-पार अंतरापृष्ठ अपनी समतल स्थिति से $\eta = sy$, $-L/2 \le y
\le L/2$ द्वारा विस्थापित होता है। घनत्व-अंतर घटी हुई परत, जो
$\eta$ से विस्थापित होती है, प्रति इकाई क्षेत्रफल $\tfrac12\rho g'\eta^2$ संचित करती
है; पट्टी पर माध्य करने से, $\langle\eta^2\rangle = s^2L^2/12$। (ख)~$g' = 0.327$:
\num{1.5e5} और \qty{6.1e5}{J/m^2}, जबकि गतिज ऊर्जा $\tfrac12\times0.7\times15^2
\times10^4 = \qty{7.9e5}{J/m^2}$ है: तुलनीय।
\end{solution}

\begin{solution}{pb:b3:atmospheric-dynamics:1}
\textbf{1.} $f_{45} = \qty{1.031e-4}{s^{-1}}$, $f_{50} =
\qty{1.117e-4}{s^{-1}}$; $\zeta = f_{45} - f_{50} = \qty{-8.6e-6}{s^{-1}}$।

\textbf{2.} $R = V/|\zeta| = 15/\num{8.6e-6} = \qty{1750}{km}$; घड़ी की सुई की दिशा में,
वापस दक्षिण की ओर।

\textbf{3.} $\zeta = f_{45}(9.5/10) - f_{50} = \qty{-1.38e-5}{s^{-1}}$, अधिक
प्रतिचक्रवातीय।

\textbf{4.} $\zeta = f_{45} - f_{40} = \qty{+9.4e-6}{s^{-1}}$, चक्रवातीय: यह
वापस उत्तर की ओर मुड़ता है। किसी भी दिशा में विस्थापित होने पर, कण वह भ्रमिलता प्राप्त करता है
जो उसे वापस मोड़ती है: एक प्रत्यानयन प्रभाव, और एक दोलन।

\textbf{5.} $\beta = \qty{1.62e-11}{m^{-1}.s^{-1}}$; वृत्त की लंबाई
\qty{28306}{km} है।

\textbf{6.} $n = 3$ से 8: तरंगदैर्घ्य 9435, 7076, 5661, 4718, 4044 और
\qty{3538}{km}; $c = -21.5$, $-5.5$, 1.9, 5.9, 8.3 और \qty{9.9}{m/s}; $c_g
= 51.5$, 35.5, 28.1, 24.1, 21.7 और \qty{20.1}{m/s}।

\textbf{7.} $n = 3$ और 4 पश्चिम की ओर चलते हैं, $n \ge 5$ पूर्व की ओर। तरंग $n = 5$ को
$10^6/1.86 = \qty{5.4e5}{s}$, \num{6.2}~दिन, \qty{1000}{km} चलने में लगते हैं।

\textbf{8.} नीचे की ओर, पूर्व: $c_g = \qty{24.1}{m/s}$ पर ऊर्जा
\qty{5000}{km} तक \num{2.4}~दिनों में फैलती है।

\textbf{9.} $\lambda_s = 2\pi\sqrt{15/\num{1.62e-11}} = \qty{6050}{km}$,
$n_s = 4.7$।

\textbf{10.} $n_s = 5.7$ \qty{10}{m/s} के लिए और 4.1 \qty{20}{m/s} के लिए: अधिक प्रबल पछुआ
लंबी स्थिर तरंगों के पक्ष में होती है, और
गोलार्ध के चारों ओर कम तरंगों के पक्ष में।

\textbf{11.} $n_s = 5.8$ \ang{35} N पर और 3.4 \ang{55} N पर।

\textbf{12.} $\zeta = f(8/10) - f = \qty{-2.06e-5}{s^{-1}}$: पर्वतों के ऊपर संपीडन
स्तंभ को प्रतिचक्रवातीय बनाता है, और पीछे की ओर विस्तार
चक्रवातीय बनाता है। पर्वत प्रवाह पर भ्रमिलता का एक स्थिर प्रारूप आरोपित करते हैं,
यानी एक ऐसा बल जो नहीं चलता।

\textbf{13.} स्थिर बल एक स्थिर अनुक्रिया उत्पन्न करता है; अनुक्रिया
सबसे प्रबल तब होती है जब आरोपित प्रारूप स्वयं एक मुक्त तरंग हो जो
स्थिर रहती, अर्थात् $\lambda_s$ पर अनुनाद।

\textbf{14.} $\Lambda = 25/10^4 = \qty{2.5e-3}{s^{-1}}$;
$|\partial T/\partial y| = \Lambda fT/g = \num{2.5e-3}\times\num{1.031e-4}
\times250/9.81 = \qty{6.6}{K}$ प्रति \qty{1000}{km}।

\textbf{15.} $\mathrm{Ri} = 0.011^2/(\num{2.5e-3})^2 = 19$: 0.25 से कहीं अधिक,
अपरूपण प्रक्षोभ नहीं।

\textbf{16.} $L_d = 0.011\times10^4/\num{1.031e-4} = \qty{1067}{km}$।

\textbf{17.} $2\pi\times1067/1.61 = \qty{4160}{km}$; $28306/4160 = 6.8$
तरंगें।

\textbf{18.} $2\pi\times1067/2.40 = \qty{2790}{km}$।

\textbf{19.} $\sigma = 0.31\times\num{1.031e-4}\times\num{2.5e-3}/0.011 =
\qty{7.3e-6}{s^{-1}}$, \num{0.63}~दिन$^{-1}$।

\textbf{20.} $\exp(3\times0.627) = 6.6$।

\textbf{21.} तूफ़ान, \qty{4160}{km}, स्थिर
तरंगदैर्घ्य \qty{6050}{km} से छोटा है: रॉस्बी तरंग के रूप में यह पछुआ द्वारा
पूर्व की ओर बहाया जाता है, और तूफ़ान धीमी या स्थिर लंबी
तरंगों के बीच से पूर्व की ओर चलते हैं।

\textbf{22.} ऊँचाई के साथ पश्चिम की ओर: केवल तभी ध्रुव की ओर गति करने वाली वायु
भूमध्य रेखा की ओर गति करने वाली वायु से अधिक गर्म होती है, जिससे ऊष्मा ध्रुव की ओर जाती है और
\omterm{def:b3:atmospheric-dynamics:ape}{उपलब्ध स्थितिज ऊर्जा} मुक्त होती है।

\textbf{23.} सबसे तेज़ बढ़ने वाला विक्षोभ हर
\textbf{\num[text-series-to-math=true]{1.6}~दिन} में $e$ गुना बढ़ता है (\num{1.59}~दिन
गणना के अनुसार)।
\end{solution}
